\begin{table}[t]
\centering
\caption{Geochemical marker suite for Module~3. Attribution relies on combinations and ratios of markers because each metal also has non-e-waste sources.}
\label{tab:markers}
\footnotesize
\begin{tabularx}{\textwidth}{@{}L{2.6cm}YL{2.55cm}Y@{}}
\toprule
Marker & Principal e-waste source & Method & Diagnostic value and caveats\\
\midrule
Cu & Wiring, printed circuit boards & ICP-MS & High loading; also agricultural and industrial sources\\
Pb, Sn & Solder; Pb also in cathode-ray-tube glass & ICP-MS & Pb--Sn covariation indicates solder; Pb also from legacy gasoline\\
Sb & Sb\textsubscript{2}O\textsubscript{3} synergist in flame-retarded plastics & ICP-MS & Tracer of the plastic fraction; also from brake wear\\
Zn, Ni, Cd & Galvanised parts, batteries & ICP-MS & Supporting markers; Cd partitioning is salinity-sensitive\\
PBDEs & Flame retardants in casings and boards & GC--MS & Highly specific; hydrophobic and particle-bound\\
PCDD/Fs, PBDD/Fs & By-products of open burning & HRGC--HRMS & Burning-specific; analysed on a targeted subset\\
\textsuperscript{206}Pb/\textsuperscript{207}Pb, \textsuperscript{208}Pb/\textsuperscript{206}Pb & Source-specific isotopic signatures & MC-ICP-MS & End-member discrimination in mixing models (\cref{eq:mix})\\
\textsuperscript{210}Pb, \textsuperscript{137}Cs & -- & Gamma spectrometry & Sediment-core chronology\\
\bottomrule
\end{tabularx}
\end{table}
